\documentclass[11pt,a4paper]{article}
\usepackage{amsmath,amssymb}
\usepackage{booktabs}
\usepackage{array}
\usepackage{xcolor}
\usepackage{listings}
\usepackage{geometry}
\usepackage{hyperref}
\usepackage{parskip}

\geometry{margin=2.5cm}

\definecolor{codebg}{rgb}{0.96,0.96,0.96}
\definecolor{codekey}{rgb}{0.0,0.4,0.7}
\definecolor{codestr}{rgb}{0.6,0.1,0.1}
\definecolor{codecom}{rgb}{0.3,0.5,0.3}

\lstset{
  language=Python,
  basicstyle=\ttfamily\small,
  keywordstyle=\color{codekey}\bfseries,
  stringstyle=\color{codestr},
  commentstyle=\color{codecom}\itshape,
  backgroundcolor=\color{codebg},
  frame=single,
  framesep=4pt,
  breaklines=true,
  numbers=none,
  showstringspaces=false,
}

\title{\textbf{Controller Evaluation Metrics}\\[4pt]
  \large Equations, definitions, and code correspondence}
\author{ROS2 HIL Visual Servoing Project}
\date{2026}

\begin{document}
\maketitle
\tableofcontents
\bigskip

% ─────────────────────────────────────────────────────────────────────────────
\section{Setup and notation}

All metrics are computed by \texttt{benchmarks/plot\_controller\_hil.py}, which
reads a recorded rosbag and writes a \texttt{metrics.csv} file.  Nothing is
measured live --- the drone flies once, then the bag is analysed offline.

\subsection*{Coordinate frame}

The simulator publishes ground-truth positions on two topics:
\begin{itemize}
  \item \texttt{/sim/drone\_pose} --- drone position $(x_d, y_d, z_d)$ at each timestep.
  \item \texttt{/sim/target\_pose} --- target (sign) position $(x_T, y_T, z_T)$,
        constant throughout a run.
\end{itemize}

\subsection*{Time origin}

Every metric is measured \emph{from the moment the benchmark FSM enters the
\texttt{APPROACHING} state} --- i.e.\ from the instant the controller takes
over after the search phase.  This time is called $t = 0$.  The full run
lasts until $t = T$ (the bag ends or the FSM transitions to \texttt{REACHED}).

% ─────────────────────────────────────────────────────────────────────────────
\section{Target standoff distance \boldmath$d^*$}

The regulation target is not zero distance --- the drone should stop at a
specific standoff, not crash into the sign.  $d^*$ is derived from the
camera's perspective model: when the target bounding box spans a fraction
$\rho^*$ of the image height, the drone is exactly at standoff.

\[
  \boxed{d^* = \frac{f_y \, H}{\rho^* \, H_{\mathrm{img}}}}
\]

\begin{center}
  \begin{tabular}{lll}
    \toprule
    Symbol & Code constant & Value \\
    \midrule
    $f_y$ & \texttt{FY} & 554.0 px (focal length, vertical) \\
    $H$ & \texttt{KNOWN\_H} & 1.5 m (physical target height) \\
    $\rho^*$ & \texttt{TARGET\_BBOX\_RATIO} & 0.55 (fraction of image height at stop) \\
    $H_{\mathrm{img}}$ & \texttt{IMG\_H} & 480 px (image height) \\
    $d^*$ & \texttt{STANDOFF} & $\approx 3.15$ m \\
    \bottomrule
  \end{tabular}
\end{center}

This comes directly from the pinhole camera projection:
a target of height $H$ at distance $d$ projects to
$h_{\mathrm{px}} = f_y H / d$ pixels.
Setting $h_{\mathrm{px}} = \rho^* H_{\mathrm{img}}$ and solving for $d$ gives the formula above.

% ─────────────────────────────────────────────────────────────────────────────
\section{Distance and error signal}

At each timestep, the 3-D Euclidean distance from the drone to the target is

\[
  d(t) = \bigl\| \mathbf{p}_T - \mathbf{p}_d(t) \bigr\|_2
        = \sqrt{(x_T - x_d)^2 + (y_T - y_d)^2 + (z_T - z_d)^2}
\]

The \emph{regulation error} --- how far the drone is from the desired standoff
--- is

\[
  e(t) = d(t) - d^*
\]

$e > 0$: drone is too far away.
$e < 0$: drone has overshot past the standoff (closer than desired).
$e = 0$: perfect standoff.

\begin{lstlisting}
dist3d = np.sqrt((tx - xs)**2 + (ty - ys)**2 + (tz - zs)**2)
err = dist3d - STANDOFF
\end{lstlisting}

% ─────────────────────────────────────────────────────────────────────────────
\section{Integral error metrics}

These three metrics summarise total accumulated error over the whole run.
They are computed by numerical quadrature (trapezoidal rule) over the
post-engage trajectory.

\subsection{IAE --- Integral of Absolute Error}

\[
  \mathrm{IAE} = \int_0^T |e(t)|\, dt
\]

The most natural summary of total error.  Units: m$\cdot$s.
A fast controller with small steady-state error gets a low IAE.

\subsection{ISE --- Integral of Squared Error}

\[
  \mathrm{ISE} = \int_0^T e(t)^2\, dt
\]

Squaring penalises large instantaneous errors more heavily than small ones.
Units: m$^2\cdot$s.
A controller that mostly stays close to $d^*$ but occasionally spikes will have
a higher ISE relative to its IAE than a controller with a smooth, uniform
approach.

\subsection{ITAE --- Integral of Time-weighted Absolute Error}

\[
  \mathrm{ITAE} = \int_0^T t\,|e(t)|\, dt
\]

The time-weighting means errors that persist \emph{late} in the run are
penalised much more than early errors (when the drone is still approaching).
Units: m$\cdot$s$^2$.
ITAE rewards controllers that settle quickly and stay settled.

\begin{lstlisting}
_trapz = np.trapezoid   # or np.trapz for older numpy
iae  = _trapz(np.abs(err), tt)
ise  = _trapz(err**2,      tt)
itae = _trapz(np.abs(err) * tt, tt)
\end{lstlisting}

\subsection*{Comparison}

\begin{center}
  \begin{tabular}{lccc}
    \toprule
    & IAE & ISE & ITAE \\
    \midrule
    Penalises large spikes? & linear & quadratic & linear \\
    Penalises late errors?  & equally & equally & more\\
    Units                   & m$\cdot$s & m$^2\cdot$s & m$\cdot$s$^2$ \\
    \bottomrule
  \end{tabular}
\end{center}

The comparison tables in the report use \textbf{IAE} as the primary integral
metric because it is the most interpretable: ``the drone accumulated $X$
metre-seconds of distance error.''

% ─────────────────────────────────────────────────────────────────────────────
\section{Settling time \boldmath$t_s$}

Settling time is the first instant at which the error enters a tolerance band
and \emph{stays there} for a hold period $T_{\mathrm{hold}}$.

\[
  \boxed{t_s = \min\bigl\{\, t \in [0,T] \;:\; |e(\tau)| \le \delta \;\; \forall\, \tau \in [t,\; t + T_{\mathrm{hold}}] \bigr\}}
\]

\begin{center}
  \begin{tabular}{lll}
    \toprule
    Symbol & Code constant & Value \\
    \midrule
    $\delta$ & \texttt{SETTLE\_BAND} & 0.30 m \\
    $T_{\mathrm{hold}}$ & \texttt{SETTLE\_HOLD} & 2.0 s \\
    \bottomrule
  \end{tabular}
\end{center}

The 2-second hold requirement filters out momentary dips through the band
that are followed by further oscillation.
If the drone never settles (e.g.\ it oscillates past $\delta$ indefinitely),
$t_s$ is recorded as \texttt{nan}.

\begin{lstlisting}
SETTLE_BAND = 0.30   # metres
SETTLE_HOLD = 2.0    # seconds

settle = float('nan')
for i in range(len(tt)):
    if abs(err[i]) <= SETTLE_BAND:
        j = i
        while j < len(tt) and tt[j] - tt[i] < SETTLE_HOLD:
            if abs(err[j]) > SETTLE_BAND:
                break
            j += 1
        else:
            settle = tt[i]; break
\end{lstlisting}

% ─────────────────────────────────────────────────────────────────────────────
\section{Steady-state error \boldmath$e_{ss}$}

Once the drone has (presumably) settled, we measure how accurately it holds
the standoff.  The last 3 seconds of the run are used as the ``steady-state
window'':

\[
  e_{ss} = \frac{1}{|W|} \int_{T-3}^{T} |e(t)|\, dt
         \approx \mathrm{mean}\bigl(|e(t)|\bigr),\quad t \in [T-3,\, T]
\]

Units: metres.  A perfect hover at $d^*$ gives $e_{ss} = 0$.

\begin{lstlisting}
tail = err[tt >= (tt[-1] - 3.0)]
ss_err = float(np.mean(np.abs(tail)))
\end{lstlisting}

% ─────────────────────────────────────────────────────────────────────────────
\section{Overshoot}

Overshoot is how far \emph{past} the standoff the drone lunged before the
controller pulled it back.

\[
  \Delta_{\mathrm{os}} = \max\!\bigl(0,\; d^* - \min_{t} d(t)\bigr)
\]

If the drone never crosses $d^*$ (it approaches and stops without
overshooting), $\Delta_{\mathrm{os}} = 0$.  Units: metres.

\begin{lstlisting}
overshoot = float(max(0.0, STANDOFF - float(np.min(dist3d))))
\end{lstlisting}

% ─────────────────────────────────────────────────────────────────────────────
\section{Path metrics}

\subsection{Path length \boldmath$L$}

The total 3-D arc length of the drone's trajectory from $t=0$ to $t=T$:

\[
  L = \sum_{k=1}^{N-1}
      \sqrt{(\Delta x_k)^2 + (\Delta y_k)^2 + (\Delta z_k)^2}
\]

where $(\Delta x_k, \Delta y_k, \Delta z_k) = \mathbf{p}_d(t_{k}) - \mathbf{p}_d(t_{k-1})$.
Units: metres.

\subsection{Path efficiency \boldmath$\eta$}

How straight was the approach, as a ratio of ideal to actual travel:

\[
  \eta = \frac{d(0) - d^*}{L}
\]

The numerator $d(0) - d^*$ is the straight-line distance the drone
\emph{needed} to cover (from its initial distance to the standoff ring).
$L$ is what it actually flew.

\begin{itemize}
  \item $\eta = 1.0$: perfectly straight beeline to the standoff ring.
  \item $\eta < 1.0$: path was longer than necessary (lateral wander, overshooting, circling).
  \item $\eta > 1.0$: impossible in practice (would require the drone to appear
        closer than $d^*$ at $t=0$, which the FSM design prevents).
\end{itemize}

\begin{lstlisting}
seg      = np.sqrt(np.diff(xs)**2 + np.diff(ys)**2 + np.diff(zs)**2)
path_len = float(np.sum(seg))
straight = float(max(0.0, dist3d[0] - STANDOFF))
path_eff = (straight / path_len) if path_len > 1e-6 else float('nan')
\end{lstlisting}

% ─────────────────────────────────────────────────────────────────────────────
\section{Control effort}

These metrics are computed from the \texttt{/cmd\_vel} topic
(\texttt{geometry\_msgs/Twist}), using the linear velocity components
$(v_x, v_y, v_z)$.

\subsection{RMS command speed}

The root-mean-square of the commanded linear speed over the run:

\[
  v_{\mathrm{rms}} = \sqrt{\frac{1}{N} \sum_{k=1}^{N} \bigl(v_{x,k}^2 + v_{y,k}^2 + v_{z,k}^2\bigr)}
\]

Units: m/s.  High $v_{\mathrm{rms}}$ means the controller commanded aggressive
motion throughout.  The velocity clamp is 3.0 m/s, so $v_{\mathrm{rms}} = 3.0$
would mean the controller was saturated for the entire run.

\subsection{Total variation (control jerk)}

Total variation measures how much the velocity command \emph{changes} between
consecutive timesteps --- a proxy for jerk:

\[
  \mathrm{TV} = \sum_{k=1}^{N-1}
      \sqrt{(\Delta v_{x,k})^2 + (\Delta v_{y,k})^2 + (\Delta v_{z,k})^2}
\]

where $\Delta v_{x,k} = v_{x,k} - v_{x,k-1}$, etc.
Units: m/s (summed velocity increments).  A smooth controller produces small
step-to-step changes; a chattery controller produces large ones.

\begin{lstlisting}
v   = np.sqrt(c[:,1]**2 + c[:,2]**2 + c[:,3]**2)
rms_v = float(np.sqrt(np.mean(v**2)))

dv       = np.sqrt(np.sum(np.diff(c[:,1:4], axis=0)**2, axis=1))
total_var = float(np.sum(dv))
\end{lstlisting}

% ─────────────────────────────────────────────────────────────────────────────
\section{Resource metrics}

These are sampled by the benchmark harness into a \texttt{cpu.csv} file
alongside the rosbag, then averaged.

\subsection{CPU usage}

\[
  \overline{\mathrm{CPU}} = \frac{1}{N} \sum_{k=1}^{N} p_k
\]

where $p_k$ is the CPU percentage (\texttt{pcpu} column) at sample $k$.
Expressed as \% of one core (so 100 = one full core, 400 = all four
RPi5 cores fully used by this process).

\subsection{Memory (RSS)}

\[
  \overline{M} = \frac{1}{N} \sum_{k=1}^{N} \frac{r_k}{1024}
\]

where $r_k$ is the resident set size in KB (\texttt{rss\_kb} column).
Result is in MB.

% ─────────────────────────────────────────────────────────────────────────────
\section{Summary table}

\begin{center}
  \begin{tabular}{llll}
    \toprule
    Metric & Symbol & Units & Ideal \\
    \midrule
    Standoff target         & $d^*$                   & m                & fixed at $\approx$3.15 m \\
    Settling time           & $t_s$                   & s                & small \\
    Steady-state error      & $e_{ss}$                & m                & small \\
    Overshoot               & $\Delta_\mathrm{os}$    & m                & 0 \\
    IAE                     & $\mathrm{IAE}$          & m$\cdot$s        & small \\
    ISE                     & $\mathrm{ISE}$          & m$^2\cdot$s      & small \\
    ITAE                    & $\mathrm{ITAE}$         & m$\cdot$s$^2$    & small \\
    Path length             & $L$                     & m                & small \\
    Path efficiency         & $\eta$                  & dimensionless    & close to 1 \\
    RMS command speed       & $v_\mathrm{rms}$        & m/s              & task-dependent \\
    Total variation (jerk)  & $\mathrm{TV}$           & m/s              & small \\
    CPU usage               & $\overline{\mathrm{CPU}}$ & \% of one core & small \\
    Memory RSS              & $\overline{M}$          & MB               & small \\
    \bottomrule
  \end{tabular}
\end{center}

% ─────────────────────────────────────────────────────────────────────────────
\section{Notes on interpretation}

\paragraph{IAE vs.\ settling time.}
These are not the same thing.  A controller can settle early but at a large
steady-state error (low $t_s$, high $e_{ss}$, low IAE if settled quickly).
Conversely, a controller that creeps slowly toward the target will have a
high IAE even if it never overshoots.

\paragraph{RMS speed and saturation.}
IBVS shows $v_\mathrm{rms} \approx 2.38$ m/s, which is 79\% of the 3.0 m/s
clamp.  This means IBVS was riding the velocity ceiling for a substantial
fraction of the approach --- explaining both its fast settling time and its
high total variation (jerk from repeatedly hitting and releasing the clamp).

\paragraph{Path efficiency and h\_vs.}
The \texttt{h\_vs} controller shows lower path efficiency than expected.
This is a structural property of the affine homography controller: the
yaw and lateral-translation error signals are degenerate (they reduce to the
same pixel-centroid term), so the controller resolves yaw before it can
translate laterally --- causing the drone to arc around rather than approach
straight.  See \texttt{docs/homography\_explained.tex}, Section~9 for the
derivation.

\paragraph{Settling band width.}
The $\pm 0.30$ m band is about 10\% of $d^* = 3.15$ m, which corresponds
roughly to a $\pm$10\% bounding-box height variation --- visible but
acceptable for outdoor signage approach.

\end{document}
